% \chapter{Phase A, Mission Design and
% Methodology}\label{chapter-6-phase-a-mission-design-and-methodology}

% Observation, this chapter needs to be rethought.

% \begin{quote}
% \textbf{Writing Guide \& Methodological Focus:}

% \begin{itemize}
% \item
%   \textbf{Objective:} To make the transition from Pre-Phase A, defining
%   the "problem space" to the Phase A "solution space" for design,
%   following a structured process.
% \item
%   \textbf{Method:} Blend a systems engineering and framework with a
%   clear description of the simulation setup.
% \item
%   \textbf{Checkpoint:} Is it perfectly clear what the thesis is trying
%   to achieve (concepts - requirements) and how the orbital mission
%   simulation is setup to find the best architecture?
% \end{itemize}
% \end{quote}

% \section{A Systems Engineering
% Approach}\label{61-a-systems-engineering-approach}

% \begin{itemize}
% \item
%   State that you will follow a structured design process analogous to
%   standard space mission design phases (e.g., based on NASA/ECSS
%   standards).
% \item
%   Outline the flow: Mission Objectives -\textgreater{} System
%   Requirements -\textgreater{} Candidate Concepts -\textgreater{} Trade
%   Space Analysis -\textgreater{} Baseline Design.
% \end{itemize}

% \section{Simulation and Analysis
% Approach}\label{62-simulation-and-analysis-approach}

% \begin{itemize}
% \item
%   \textbf{6.3.1 Trade Space Definition:}
% \item
%   \textbf{6.3.2 Simulation Environment and Tools:}
% \item
%   \textbf{6.3.3 Analysis Method:}
% \end{itemize}
